En este trabajo se realiza un análisis experimental del rendimiento práctico de algoritmos de ordenamiento (\textsc{sort}, \textsc{QuickSort}, \textsc{MergeSort} y \textsc{PatienceSort}) y de multiplicación de matrices cuadradas (\textsc{Naive} y \textsc{Strassen}), contrastando sus tiempos de ejecución y consumo de memoria con su complejidad asintótica teórica. Evaluados bajo diversas dimensiones, distribuciones iniciales y dominios numéricos, los resultados confirman la validez asintótica teórica y evidencian cómo las constantes ocultas, determinan la eficiencia real de los algoritmos.